
\section{Fixed Workflow and Control Patterns}
\subsection{Sequential prompt chaining}
The output of one defined stage becomes input to the next; validators run between stages. Appropriate for extraction--normalization--drafting--verification, compliance form pipelines, and ordered document production. Failure of a stage should stop or move to a typed recovery branch.

\begin{figure}[htbp]\centering
\resizebox{0.96\linewidth}{!}{%
\begin{tikzpicture}
\node[flowbox] (in) at (0,0) {Typed input};
\node[flowbox] (a) at (3,0) {Model A\\Extract};
\node[flowcontrol] (gate) at (6,0) {Validate output};
\node[flowbox] (b) at (9,0) {Model B\\Transform};
\node[flowcontrol] (out) at (12,0) {Final verification};
\foreach \a/\b in {in/a,a/gate,gate/b,b/out} {\draw[flowline] (\a)--(\b);}
\end{tikzpicture}}
\caption{Prompt chaining with a deterministic intermediate gate.}
\label{fig:chaining}
\end{figure}

\textbf{Strengths:} explicit contracts, repeatability of orchestration, simple tests. \textbf{Costs:} summed stage latency and error propagation. \textbf{Controls:} per-stage validation, bounded retry on transient failure, versioned transformations, dead-letter outcomes. Approximate critical path: $T_{\mathrm{chain}}\simeq\sum_{i=1}^{n}(T_i+V_i)$, where $V_i$ represents validation or handoff cost.

\subsection{Parallelization: sectioning and voting}
Sectioning assigns disjoint independent tasks to workers. Voting obtains multiple perspectives or attempts for the same task. Voting does \emph{not} create an independent ground-truth oracle when the models share data or failure modes. Aggregation must handle missing, late, duplicate, and conflicting results.

\begin{figure}[htbp]\centering
\resizebox{0.94\linewidth}{!}{%
\begin{tikzpicture}
\node[flowbox] (in) at (0,0) {Input};
\node[flowbox] (a) at (3,1.4) {Worker A\\Section one};
\node[flowbox] (b) at (3,0) {Worker B\\Section two};
\node[flowbox] (c) at (3,-1.4) {Worker C\\Section three};
\node[flowcontrol] (join) at (6.5,0) {Join and reconcile};
\node[flowbox] (out) at (10,0) {Validated result};
\foreach \w in {a,b,c} {\draw[flowline] (in.east)--(\w.west); \draw[flowline] (\w.east)--(join.west);}
\draw[flowline] (join)--(out);
\end{tikzpicture}}
\caption{Static fan-out/fan-in. Parallel workers do not automatically imply multi-agent autonomy.}
\label{fig:parallel}
\end{figure}

\textbf{Strengths:} shorter critical path for independent tasks, specialized checks. \textbf{Costs:} concurrency, duplicate tokens, stragglers, integration failures. \textbf{Controls:} bounded fan-out, worker deadlines, explicit all-required/quorum policies, deterministic conflict handling. Approximation: $T_{\mathrm{parallel}}\simeq T_{\mathrm{dispatch}}+\max(T_i)+T_{\mathrm{join}}$.

\subsection{Routing and dispatch}
A classifier selects a specialized path based on validated intent, constraints, and risk. Implement confidence/unknown handling as a typed result; use a safe fallback rather than guessing at an unauthorized action. Routing decisions can be deterministic rules, a small classifier, or a model, but downstream execution permissions remain independent.

\begin{figure}[htbp]\centering
\resizebox{0.95\linewidth}{!}{%
\begin{tikzpicture}
\node[flowbox] (in) at (0,0) {Request};
\node[flowcontrol] (route) at (3,0) {Schema and intent router};
\node[flowbox] (a) at (6,1.3) {Read-only research};
\node[flowbox] (b) at (6,0) {Code assessment};
\node[flowreview] (c) at (6,-1.3) {Unknown or high risk};
\node[flowcontrol] (norm) at (9.6,0) {Normalize and verify};
\draw[flowline] (in)--(route);
\foreach \w in {a,b,c} {\draw[flowline] (route.east)--(\w.west); \draw[flowline] (\w.east)--(norm.west);}
\end{tikzpicture}}
\caption{Routing with a first-class fallback for ambiguous or high-risk requests.}
\label{fig:routing}
\end{figure}

\textbf{Strengths:} lower average cost, specialist prompts and tools, clear isolation. \textbf{Costs:} misclassification, prompt-injection-induced routing errors, changing categories. \textbf{Controls:} golden-set intent tests, explicit uncertainty, deny-by-default mutation paths, audit of route decisions.

\subsection{Deterministic control enclosing agentic reasoning}
The model may interpret evidence or draft an action; code implements valid transitions and executes only authorized changes. This approach is essential whenever reproducible policy, legal/regulatory controls, or expensive irreversible outcomes dominate the requirement \cite{databricks}. Determinism here concerns \emph{the authority decision}, not the phrasing of model-generated output.

\begin{notebox}[Common composition]
A regulated process may combine a router, parallel read-only investigation, an evaluator, and a human approval checkpoint while its outer state machine remains deterministic.
\end{notebox}
